\section{Boundary-uniform coding of preparation instruments}\label{sec:boundarycode68}
Strict positivity is essential to the two-point programme argument, as
Proposition~\ref{prop:boundaryphase67} shows. Nevertheless it is not
necessary for a complete coding law in an irreversible family. We
consider all outcome-retaining instruments that erase their input and
prepare a specified quantum--classical state. This family contains zero
outcomes and arbitrary rank-deficient output blocks.

For $m,n\ge1$ put
\[
 \mathcal S_{m,n}=\{(\sigma_y)_{y=1}^m:\sigma_y\in\operatorname{Herm}(n),
                  \ \sigma_y\ge0,\ \sum_y\tr\sigma_y=1\}.
\]
Write $\boldsymbol\sigma=\bigoplus_y\sigma_y$. For any input
dimension $d\ge1$, the associated instrument is
\begin{equation}\label{eq:replacer68}
 \mathcal R_{\sigma,y}(X)=\tr(X)\sigma_y,\qquad
 J_y(\mathcal R_\sigma)=I_d\otimes\sigma_y.
\end{equation}
The output label $y$ is retained. Its conditional quantum state is
$\sigma_y/\tr\sigma_y$ when that trace is positive. This is a
specified mathematical quantum instrument; a description of it is not
a numerical algorithm acting on an unknown physical input.

\begin{lemma}[Reduction to a product-state experiment]\label{lem:productreplacer68}
For arbitrary $\sigma,\tau\in\mathcal S_{m,n}$,
\begin{equation}\label{eq:productreplacer68}
 d_N(\mathcal R_\sigma,\mathcal R_\tau)
   =\|\boldsymbol\sigma^{\otimes N}
                -\boldsymbol\tau^{\otimes N}\|_1.
\end{equation}
Here the tester may have arbitrary quantum memory, initial reference
entanglement, feedback and a public stopping time at most $N$.
\end{lemma}
\begin{proof}
Supply $N$ independent program states $\boldsymbol\sigma$ to one
fixed processor. At the next device call the processor discards the
incoming device system and returns the next program state, including
its classical outcome. Between calls it performs the prescribed tester
operations. This reproduces the experiment exactly, even when the
incoming system is entangled with the tester memory: the incoming
system is traced out in \eqref{eq:replacer68}. Unused program states
are ignored after a public stop. The processor is the same channel
for $\sigma$ and $\tau$, so trace-norm contraction gives the upper
bound. Conversely, a tester can make exactly $N$ calls and retain
every outcome and quantum output. Its final state is the displayed
product, which proves equality. The processor is a quantum channel,
not a claimed finite-classical-message physical simulator.
\end{proof}

\subsection{Triangular factors and a rational code}
Every $\sigma\in\mathcal S_{m,n}$ has triangular factors
$L_y$ with $\sigma_y=L_yL_y^*$ and nonnegative real diagonal entries.
The usual no-pivot Cholesky procedure also covers singular matrices:
a zero diagonal pivot forces the entire residual row and column to
vanish and contributes a zero factor column. A block of rank $r$
therefore has exactly $r$ positive pivots. All factors together satisfy
$\sum_y\|L_y\|_F^2=1$.

Fix public bounds $0\le r_y\le n$, not all zero, and write
\begin{equation}\label{eq:rankdimension68}
 D=\sum_{y=1}^m r_y(2n-r_y),\qquad v=D-1,
 \qquad
 \mathcal S_{\boldsymbol r}
   =\{\sigma\in\mathcal S_{m,n}:\operatorname{rank}\sigma_y\le r_y\}.
\end{equation}
If the positive pivot positions of $L_y$ are $T_y\subseteq\{1,\ldots,n\}$,
record only the real diagonal and the real and imaginary coordinates
strictly below each such pivot. Their total number is
\[
 h=\sum_y\sum_{j\in T_y}\{1+2(n-j)\}\le D.
\]
The resulting real coordinate vector $x\in\R^h$ has norm one.
The inequality follows by placing each block's pivots as early as
possible. It holds also when the actual ranks are smaller than the
public bounds.

Choose a coordinate $j$ of maximal absolute value, put $a=|x_j|$,
and set $w=x/a$. For an integer $B\ge1$, take
\begin{equation}\label{eq:factorgrid68}
 z_j=B\operatorname{sgn}(x_j),\qquad
 z_i=\operatorname{sgn}(x_i)\lfloor B|w_i|\rfloor\quad(i\ne j).
\end{equation}
Insert these integers in the recorded factor positions to obtain
Gaussian-integer triangular matrices $Z_y$; all unrecorded columns
are zero. Define
\begin{equation}\label{eq:factoroutput68}
 T=\sum_y\|Z_y\|_F^2,\qquad
 \widehat\sigma_y=Z_yZ_y^*/T.
\end{equation}
Since $|z_j|=B$, the denominator satisfies $B^2\le T\le hB^2$.

\begin{theorem}[A boundary-uniform rational factor code]\label{thm:factorcode68}
The output \eqref{eq:factoroutput68} lies in
$\mathcal S_{\boldsymbol r}$, preserves every zero target outcome,
and does not increase the rank of any individual block. For all
$N\ge1$,
\begin{equation}\label{eq:factorerror68}
 d_N(\mathcal R_\sigma,\mathcal R_{\widehat\sigma})
 \le\min\left\{2,\frac{4\sqrt{N(h-1)}}B\right\}
 \le\min\left\{2,\frac{4\sqrt{Nv}}B\right\}.
\end{equation}
At fixed $m,n,\boldsymbol r,B$ a fixed-length code needs at most
\begin{equation}\label{eq:factorbits68}
 mn+\left\lceil\log_2\bigl(2D(2B+1)^v\bigr)\right\rceil
\end{equation}
bits. For rational target blocks, its encoder and decoder use exact
rational and integer arithmetic, including integer square roots; no
algebraic-number oracle, positive eigenvalue margin, or Kraus
factorization is required.
\end{theorem}
\begin{proof}
The product $Z_yZ_y^*$ is positive semidefinite, and normalization
makes the sum of its traces exactly one. Its number of possibly nonzero
columns is no larger than that of $L_y$, so its rank does not exceed
the target rank. When an outcome is zero it has no recorded columns
and remains exactly zero.

Put $z'=z/B$. Its $j$th coordinate agrees with that of $w$; every
other coordinate differs by less than $1/B$. Both $\|w\|$ and
$\|z'\|$ are at least one. The elementary normalization estimate gives
\begin{equation}\label{eq:factornormalize68}
 \left\|x-\frac z{\|z\|}\right\|
 \le\frac{2\|w-z'\|}{\|w\|}
 \le\frac{2\sqrt{h-1}}B.
\end{equation}
Use the normalized vectors
\[
 |L\rangle=\sum_{y,i,j}(L_y)_{ij}|y,i\rangle|y,j\rangle,
 \qquad
 |\widehat L\rangle=T^{-1/2}\sum_{y,i,j}(Z_y)_{ij}|y,i\rangle|y,j\rangle.
\]
Their reduced states on the first factor are
$\boldsymbol\sigma$ and $\widehat{\boldsymbol\sigma}$. Their vector
distance is the left side of \eqref{eq:factornormalize68}. For any
unit vectors $u,v$, diagonalization on their two-dimensional span gives
\[
 \||u\rangle\langle u|^{\otimes N}
       -|v\rangle\langle v|^{\otimes N}\|_1
   =2\sqrt{1-|\langle u,v\rangle|^{2N}}
   \le2\sqrt N\,\|u-v\|.
\]
Indeed $1-t^N\le N(1-t)$ on $[0,1]$ and
$1-|\langle u,v\rangle|^2\le2(1-\Re\langle u,v\rangle)$.
Taking partial traces and using Lemma~\ref{lem:productreplacer68}
proves \eqref{eq:factorerror68}. These are standard purification and
product-fidelity estimates; see \cite[Section~3.2]{Watrous65}. Their
short proof specifies the normalization used here.

The pivot sets use $mn$ bits. Given them, there are at most $2h$
anchor/sign choices and $(2B+1)^{h-1}$ remaining digit strings.
Pad their index to the maximum length in \eqref{eq:factorbits68}.
An invalid header may be rejected, or sent by a total mathematical
decoder to a fixed legal state. Valid factor words always decode
positively, even if they did not arise from this encoder.

For rational input, maintain the exact rational Schur residual $A$.
At a positive pivot $u=A_{jj}$, the factor diagonal is $\sqrt u$
and a lower entry is $A_{ij}/\sqrt u$. Thus the square of each of
its real coordinates is rational, respectively $u$, or
$(\Re A_{ij})^2/u$, or $(\Im A_{ij})^2/u$; its sign is known.
Updating the residual uses rational arithmetic only. Choose the
largest squared coordinate $b_j>0$ by exact comparisons. If another
coordinate has squared value $b_i$, then
\[
 \lfloor B|w_i|\rfloor
   =\operatorname{isqrt}\!\left(\left\lfloor B^2b_i/b_j\right\rfloor\right).
\]
Consequently no square root of a rational number must be represented.
The ratios-of-minors bounds in Lemma~\ref{lem:schurminors67} give
polynomial operand lengths during this preprocessing. Integer square
root, products and mixed-radix coding have polynomial bit complexity
in the explicit input and output lengths. Encoding workspace and the
expanded decoded matrix are not identified with the compressed payload.
\end{proof}

\subsection{Sharp description laws on the boundary}
Let $\mathcal P_N(\delta;\boldsymbol r)$ be the minimum number of
legal preparation instruments whose $d_N$-balls of radius $\delta$
cover $\mathcal S_{\boldsymbol r}$; the decoded centre may have any
output rank. Add a superscript $\Q$ when rational centres are required.
For an upper construction our centres will obey the same rank bounds.

\begin{theorem}[Rank-stratified joint coding law]\label{thm:rankentropy68}
Fix $m,n,\boldsymbol r$ as in \eqref{eq:rankdimension68}. If $v>0$,
there exist constants $c,C>0$, depending only on these fixed data,
such that, for every $N\ge1$ and $0<\delta\le1/32$,
\begin{equation}\label{eq:rankcover68}
 \max\{1,(c\sqrt N/\delta)^v\}
 \le\mathcal P_N(\delta;\boldsymbol r)
 \le\mathcal P_N^{\Q}(\delta;\boldsymbol r)
 \le C(\sqrt N/\delta)^v.
\end{equation}
The optimal legal and rational fixed-length descriptions therefore have
length
\begin{equation}\label{eq:rankbits68}
 \frac v2\log_2N+v\log_2(1/\delta)+O_{m,n,\boldsymbol r}(1)
\end{equation}
uniformly in $N$ and $\delta$. The statement has no spectral-gap or
positive-eigenvalue-margin hypothesis. It includes every rank-deficient
boundary allowed by the public bounds. If $v=0$, the family is a
singleton and its optimal description length is zero.
\end{theorem}
\begin{proof}
For the upper bound take $B=\lceil4\sqrt{Nv}/\delta\rceil$ in
Theorem~\ref{thm:factorcode68}; for rational implementation it suffices
to replace $\sqrt{Nv}$ by its integer ceiling. Its error is at most
$\delta$, and its codebook cardinality is at most
$2^{mn}2D(2B+1)^v$. This has the asserted upper order uniformly on
the stated parameter range.

For completeness we construct a full-dimensional local patch for the
lower bound. For every block with $r_y>0$, take an $n$ by $r_y$
lower-trapezoidal factor: entries above the first $r_y$ diagonal
positions vanish, and these diagonal entries are positive real
numbers. These factors have $D$ independent real coordinates. Impose
$\sum_y\|L_y\|_F^2=1$ and work near the point with all these diagonal
entries equal to $(\sum_y r_y)^{-1/2}$ and all other entries zero.
An open patch of this unit sphere is a smooth $v$-dimensional manifold.
The map
\[
                         (L_y)_y\longmapsto(L_yL_y^*)_y
\]
is one-to-one there and has a smooth inverse on its image. To see the
inverse explicitly, the leading $r_y$ by $r_y$ principal block is
positive definite and determines its unique positive-diagonal Cholesky
factor $A_y$. The bottom-left block is $B_yA_y^*$, so it determines
$B_y$ by multiplication by $(A_y^*)^{-1}$. These operations are smooth
near the chosen point. When $r_y=0$, the entire block is identically
zero and contributes no coordinate.

Choose a closed coordinate ball of sufficiently small positive radius
inside this patch. The chart and its inverse have bounded derivatives
on a slightly larger neighborhood; hence distances between parameters
in the ball are bounded above and below by constant multiples of the
tuple-Frobenius distance between their states. More explicitly one
may bound the derivative of the inverse on a convex neighborhood of
the chosen positive principal blocks, where the same Cholesky formulas
extend smoothly. A maximal Euclidean packing of this parameter ball
at scale $A\delta/\sqrt N$ therefore has at least
$(c\sqrt N/\delta)^v$ points whose state distances exceed
$32\delta/\sqrt N$, after increasing the fixed $A$.

The binary argument of Corollary~\ref{cor:localchoi68}, applied
directly to the trace-one block states (input dimension one), separates
each pair by more than $2\delta$ after $N$ copies. By
Lemma~\ref{lem:productreplacer68} the same separation holds in $d_N$.
One radius-$\delta$ ball, regardless of the rank of its legal centre,
cannot contain two packing points. This proves the lower bound and
then \eqref{eq:rankbits68}. Finally $D=1$ requires $n=1$ and exactly
one bound $r_y=1$; all other bounds vanish. The sole permitted state
has unit mass on that outcome, as asserted.
\end{proof}

\begin{corollary}[All ranks, with an explicit lower constant]\label{cor:fullboundary68}
For the full family $\mathcal S_{m,n}$ put $v=mn^2-1>0$. Then
\begin{equation}\label{eq:fullboundary68}
 \max\left\{1,\left(\frac{\sqrt N}{64mn\delta}\right)^v\right\}
 \le\mathcal P_N(\delta)\le\mathcal P_N^{\Q}(\delta)
 \le 2^{mn}2mn^2(2B+1)^v,
 \qquad B=\left\lceil\frac{4\sqrt{Nv}}\delta\right\rceil.
\end{equation}
In particular a general $n$-dimensional prepared state has coefficient
$n^2-1$, while the class of pure prepared states has coefficient
$2n-2$ in \eqref{eq:rankbits68}. A classical $m$-outcome law has
coefficient $m-1$, including probabilities equal to zero.
\end{corollary}
\begin{proof}
The upper bound is Theorem~\ref{thm:factorcode68} with every $r_y=n$.
For the lower bound, the affine space of block states has dimension
$mn^2-1$. Its tuple-Frobenius ball of radius $1/(2mn)$ about the
blocks $I_n/(mn)$ lies in $\mathcal S_{m,n}$. Repeat the packing proof
of Theorem~\ref{thm:jointcode68} with input dimension one. For pure
states take $m=1,r_1=1$ in Theorem~\ref{thm:rankentropy68}; for
classical laws take $n=1$.
\end{proof}

\begin{remark}[Operational and priority distinctions]\label{rem:boundaryscope68}
Input erasure is decisive in \eqref{eq:productreplacer68}. No analogous
product-state reduction has been proved here for an arbitrary disturbing
instrument; coherent accumulation for unitary channels remains governed
by Proposition~\ref{prop:boundaryphase67}. The present theorems classify
reusable description lengths in the stated preparation families, not
general simulator workspace. Constants remain dimension dependent.

Quantum population compression starts with physical copies of an
unknown state and compresses their quantum information. Yang, Bai,
Chiribella and Hayashi \cite{YBCH68} obtain a half-parameter-count
logarithmic memory law in that different task and show a limitation
on purely classical compression of noncommuting families. Here the
encoder is given a classical matrix description, and each codeword
specifies a state or preparation instrument. It need not infer that
word from unknown copies. The common half-dimension scale, purification
inequalities and finite-dimensional covering methods are classical
antecedents, not claimed inventions. The contribution established here
is a uniform two-parameter legal-description law with a rational,
zero-outcome-preserving, rank-nonincreasing encoder and matching lower
bounds throughout the declared boundary families.
\end{remark}
